\subsection{Poincare Group}
\label{subsec:poincare-group}

The Poincare group is the group of isometries of Minkowski spacetime. It is a semidirect product of the translation group and the Lorentz group: $\text{ISO}(1,3) = \mathbb{R}^{1,3} \rtimes SO(1,3)$. The Poincare group is a 10-dimensional Lie group, with 4 translation generators and 6 Lorentz generators.

A particle (fundamental or composite) gives a unitary representation of the Poincare group: the transormations in the Poincare group act on the states of the particle, transforming its coordinates, while leaving the mass and spin invariant. A particle is fundamental if it gives irreducible representations of the Poincare group. Thus, we can define \textbf{particles} as irreducible unitary representations of the Poincare group.

All unitary representations of the Poincare group are infinite-dimensional (Fact~\ref{fact:infinite-dim}). The tensors, for example, fail to be unitary representations.  By Wigner's theorem, the irreducible unitary representations of the Poincare group are classified by the mass $m$ and spin $J$ of the particle. When $J > 0$, for each value of momentum with $p^2 = m^2$ there are $2 J + 1$ independent states in the representation (in which case they are called \textit{spin states}), and only $2$ states when $m = 0$ (in which case they are called \textit{helicity states}). If $J = 0$ there is only one state for each value of momentum (in which case it has spin-$0$, or \textit{spinless}, and the particle acts like a point mass).

